\section{Limitations and Future Work}
\label{sec:limitations}

PrefBench is a benchmark prototype, and its limitations are part of the release contract.

\paragraph{Simulator-based data.}
The buyer population is semi-synthetic. Observable variables are configured from plausible population assumptions, while hidden preferences and negotiation traits are generated by hand-designed conditional rules. The benchmark should therefore be used for controlled study of personalized pricing agents, not as evidence of deployable real-world pricing performance.

\paragraph{Limited external validation.}
The current release does not include a large human-subject study or real transaction data. Future work should add human reference evaluation, expert review of buyer behavior, or small-scale preference elicitation to calibrate parts of the latent model.

\paragraph{Limited LLM evaluation.}
The current prototype is designed to support future LLM-facing evaluation, but the reported results focus on heuristic baselines. A natural next step is to add a natural-language observation renderer, a structured JSON action schema, and low-cost zero-shot LLM baselines.

\paragraph{Diagnostic extensions.}
The hidden variables in the simulator make belief probing possible: one can ask whether an agent or probe model can infer price sensitivity, patience, bargaining strength, or willingness-to-pay buckets from partial negotiation history. Robustness suites and out-of-distribution persona splits are also important future extensions.
